\section{Affine-slice recurrence for \texorpdfstring{\seq{A005787}}{A005787}}
\label{sec:a005787}

Consider an arrangement of \(n\) circles and identify its \(2^n\) regions with
\(\{0,1\}^n\), where \(x_i\) records membership in circle \(i\).  A black-white
coloring is a Boolean function.  Starting from the all-white coloring, the
allowed operation on coordinate \(i\) either colors the whole circle black or
reverses every color inside it.  Algebraically these are
\begin{equation}
 f\longmapsto f\lor x_i,
 \qquad
 f\longmapsto f\oplus x_i,
 \label{eq:a005787-ops}
\end{equation}
where arithmetic for \(\oplus\) is over \(\F_2\).  This is the subset-coloring
problem of Rubinstein, Shallit, and Szegedy
\citep{RubinsteinShallitSzegedy1988}.

Let \(F_n\) be the family of reachable functions and put
\(a(n)=|F_n|\), with the auxiliary value \(a(0)=1\).  Let
\(\operatorname{Aff}_n\) be the \(2^{n+1}\) affine functions on \(n\) variables.

\subsection{The last-OR decomposition}

For \(0\le i<n\), define
\begin{equation}
 E_i=\left\{f:
 f|_{x_i=0}\in F_{n-1}
 \text{ and }
 f|_{x_i=1}\in\operatorname{Aff}_{n-1}
 \right\}.
 \label{eq:a005787-E}
\end{equation}

\begin{lemma}
\label{lem:a005787-union}
One has
\begin{equation}
 F_n=\bigcup_{i=0}^{n-1}E_i,
 \qquad
 |E_i|=2^n a(n-1).
 \label{eq:a005787-union}
\end{equation}
\end{lemma}

\begin{proof}
Choose the last OR operation in a construction.  If it is performed on
coordinate \(i\), subsequent operations are XORs by coordinate functions.
The \(x_i=0\) slice is therefore reachable in \(n-1\) variables and the
\(x_i=1\) slice is affine.  A construction using no OR is itself affine and
belongs to every \(E_i\).  Hence \(F_n\subseteq\bigcup_iE_i\).

Conversely, suppose the zero-slice is \(g\in F_{n-1}\) and the one-slice is
\(h=c\oplus\ell\), where \(\ell\) is linear in the remaining variables.
Reachability is closed under XOR by a linear form.  Construct
\(g\oplus\ell\), apply OR on coordinate \(i\), and then XOR by \(\ell\).
The two slices are now \(g\) and \(1\oplus\ell\).  If \(c=0\), one final XOR
by \(x_i\) changes only the one-slice and gives \(\ell\); if \(c=1\), no final
step is needed.  Thus every element of \(E_i\) is reachable.

There are \(a(n-1)\) choices for the zero-slice and \(2^n\) affine functions
on the remaining \(n-1\) variables, proving the cardinality statement.
\end{proof}

Restriction to any collection of zero-coordinate faces preserves
reachability: each operation restricts either to the corresponding operation in
fewer variables or to the identity.

\subsection{Multiple intersections}

\begin{lemma}[Affine patching]
\label{lem:a005787-intersection}
For every \(I\subseteq\{0,\ldots,n-1\}\) with \(|I|=k\ge2\),
\begin{equation}
 \left|\bigcap_{i\in I}E_i\right|=2^{n+1}a(n-k).
 \label{eq:a005787-intersection}
\end{equation}
\end{lemma}

\begin{proof}
Take \(f\in\bigcap_{i\in I}E_i\).  For each \(i\in I\), the slice
\(h_i=f|_{x_i=1}\) is affine, and the slices agree on every pairwise overlap.
Any two compatible affine slices glue to a global affine function: their
constant and linear coefficients agree on the common codimension-two face,
and the two remaining coordinate coefficients are fixed by the respective
slices.  Let \(h\) be the affine function obtained from two slices.  For a
third \(k\in I\), the difference
\(h|_{x_k=1}-h_k\) is affine on that slice and vanishes on its intersections
with two distinct coordinate hyperplanes.  The zero set of a nonzero affine
function over \(\F_2\) is a single affine hyperplane, so this difference is
zero.  Repeating the argument shows that a unique global affine function \(h\)
agrees with \(f\) wherever at least one selected coordinate is one.

On the remaining face
\[
 Z_I=\{x:x_i=0\text{ for every }i\in I\},
\]
put \(q=f|_{Z_I}\).  By restriction, \(q\in F_{n-k}\).  Hence \(f\) determines
a pair
\[
 (h,q)\in\operatorname{Aff}_n\times F_{n-k}.
\]

Conversely, given such a pair, define \(f=h\) off \(Z_I\) and \(f=q\) on
\(Z_I\).  We prove by induction on \(k\) that \(f\in\bigcap_{i\in I}E_i\).
Fix \(i\in I\).  The one-slice is affine.  Its zero-slice is the analogous
patching construction in one fewer variable with selected set
\(I\setminus\{i\}\).  When this set has one element, the converse part of
\cref{lem:a005787-union} shows that the slice is reachable.  For larger sets,
the induction hypothesis shows that the slice lies in the corresponding
intersection and hence is reachable.  Thus \(f\in E_i\) for every \(i\in I\).
The two constructions are inverse bijections.  Since
\(|\operatorname{Aff}_n|=2^{n+1}\), \cref{eq:a005787-intersection} follows.
\end{proof}

\subsection{Recurrence and generating function}

\begin{theorem}[All-\(n\) recurrence]
\label{thm:a005787-recurrence}
With \(a(0)=1\),
\begin{equation}
 \boxed{
 a(n)=n2^n a(n-1)+2^{n+1}
 \sum_{k=2}^{n}(-1)^{k+1}\binom{n}{k}a(n-k).}
 \label{eq:a005787-recurrence}
\end{equation}
\end{theorem}

\begin{proof}
Apply inclusion-exclusion to \(F_n=\bigcup_iE_i\).  The singleton terms use
\(|E_i|=2^na(n-1)\), and every \(k\)-fold intersection with \(k\ge2\) has the
cardinality in \cref{eq:a005787-intersection}.  There are \(\binom nk\)
intersections of size \(k\), giving \cref{eq:a005787-recurrence}.
\end{proof}

Let
\[
 A(x)=\sum_{n\ge0}a(n)\frac{x^n}{n!}.
\]
Multiplying \cref{eq:a005787-recurrence} by \(x^n/n!\), summing over \(n\),
and collecting the binomial convolution gives the formal functional equation
\begin{equation}
 \boxed{A(x)=1+2\bigl(1-x-e^{-2x}\bigr)A(2x).}
 \label{eq:a005787-egf}
\end{equation}
Set \(C(x)=2(1-x-e^{-2x})\).  Since \(C(0)=0\), iterating
\cref{eq:a005787-egf} is coefficientwise legitimate in \(\Q[[x]]\) and yields
\begin{equation}
 \boxed{
 A(x)=\sum_{m\ge0}\prod_{j=0}^{m-1}C(2^jx)
 =\sum_{m\ge0}2^m\prod_{j=0}^{m-1}
 \left(1-2^jx-e^{-2^{j+1}x}\right).}
 \label{eq:a005787-product}
\end{equation}
The empty product is one, and the \(m\)-th summand has valuation at least
\(m\), so only finitely many summands affect any coefficient.

The recurrence begins
\[
 2,\ 8,\ 112,\ 5856,\ 869824,\ 323041664,
 \ 284820663040,\ 578706325429760,\ldots.
\]
An explicit construction of all small intersection sets verifies
\cref{eq:a005787-intersection} through four variables.  A separate C++
enumerator independently obtains the complete six-circle count and its
intersection table.  Exact recurrence arithmetic through \(n=20\) agrees with
the preserved certificate.
